\documentclass{article}
\usepackage[T1]{fontenc}
\usepackage{lmodern}
\usepackage{graphicx}
\usepackage{amsmath,amssymb}
\usepackage[a4paper,margin=2cm]{geometry}
\usepackage[absolute,overlay]{textpos}
\setlength{\TPHorizModule}{1mm}
\setlength{\TPVertModule}{1mm}
\usepackage[indonesian]{babel}
\usepackage{hyperref}
\setlength{\emergencystretch}{2em}
% unit: Halaman 51

\begin{document}

% === Halaman 51 (llm) ===

\section*{Aplikasi ECC}

\begin{itemize}
    \item Banyak piranti yang berukuran kecil dan memiliki keterbatasan memori dan kemampuan pemrosesan.
    \item Di mana kita dapat menerapkan ECC?
    \begin{itemize}
        \item Piranti komunikasi nirkabel
        \item \textit{Smart cards}
        \item Web server yang membutuhkan penanganan banyak sesi enkripsi
        \item \textcolor{red}{Sembarang aplikasi yang membutuhkan keamanan tetapi memiliki kekurangan dalam \textit{power}, \textit{storage} and kemampuan komputasi adalah potensial memerlukan ECC}
    \end{itemize}
\end{itemize}

\vspace{10mm}

\centering
\small{*) Sumber bahan: \textcolor{red}{Debdeep Mukhopadhyay}, \textcolor{red}{Elliptic Curve Cryptography}, \\ \textcolor{red}{Dept of Computer Sc and Engg IIT Madras}}

\end{document}
